\documentclass[12pt,a4paper]{article}
\usepackage[utf8]{inputenc}

\setcounter{page}{1}
\pagenumbering{arabic}
\usepackage{amsmath,amsfonts,amssymb,amsthm}
\usepackage{enumerate}
\usepackage[dvips]{graphicx,epsfig}
\usepackage{setspace}
\usepackage{geometry}
\usepackage{hyperref}
\usepackage{color}
\usepackage[british]{babel}
\usepackage[mmddyyyy]{datetime}
\renewcommand{\baselinestretch}{1.2}

\marginparwidth=0in \marginparsep=0in \oddsidemargin=0in \evensidemargin=0in \headheight=0in \headsep=0in \topmargin=0.5in \textheight=9in
\textwidth=6in
\setlength{\parindent}{0pt}

\begin{document}
\begin{tabular}{l r}
University of Essex  \hspace{35ex}& Session 2025-2026 \\
Department of Economics \hspace{35ex} & Autumn Term\\
Dr Ben Etheridge \hspace{35ex} &
\end{tabular}
\\
\begin{center}
\large{\textbf{EC466: Econometrics }\\
Problem Set \# 6: Difference-in-Differences}
\end{center}
\vspace*{0.5cm}

\textbf{Background:} A government training programme for unemployed workers was implemented in Region A starting in January 2023, while Region B did not implement the programme. You have panel data tracking the same individuals in both regions from 2022 and 2023, with the following average earnings (in \pounds1000s):

\vspace{1em}
\begin{center}
\begin{tabular}{lccc}
\hline
Region & 2022 & 2023 & Sample Size ($n$) \\
\hline
A (Treated) & 18.5 & 24.2 & 250 \\
B (Control) & 20.1 & 22.8 & 300 \\
\hline
\end{tabular}
\end{center}
\vspace{1em}

\begin{enumerate}
\item \textbf{DiD Estimation and Identification}
\begin{enumerate}[(a)]
\item Calculate the canonical $2\times2$ difference-in-differences estimate of the average treatment effect on the treated (ATT). Show your working and interpret the result in the context of the training programme.
\item State the three key assumptions (SUTVA, No-Anticipation, and Parallel Trends) required for the DiD estimator to identify the ATT. Explain each assumption in the context of this application.
\item Explain intuitively why the parallel trends assumption is necessary. What would happen to the DiD estimate if this assumption were violated?
\end{enumerate}

\item \textbf{Potential Outcomes and Missing Data}\\
Consider the potential outcomes framework where $Y_{i,t}(g)$ denotes the earnings individual $i$ would have in period $t$ if they belonged to group $g \in \{A, B\}$.
\begin{enumerate}[(a)]
\item Under SUTVA and No-Anticipation, show which potential outcomes are observed for individuals in each region and time period. Present your answer in a table similar to the one on page 11 of the lecture notes.
\item Write down the formal expression for the selection bias that would arise from a simple comparison of post-treatment outcomes: $E[Y_{i,t=2023}|G_i = A] - E[Y_{i,t=2023}|G_i = B]$. Explain why this does not generally equal the ATT.
\item Using the Parallel Trends Assumption, show algebraically how the unobserved counterfactual $E[Y_{i,t=2023}(\infty)|G_i = A]$ can be imputed from observed data.
\end{enumerate}

\item \textbf{Regression Specification and Interpretation}\\
Suppose you estimate the following two-way fixed effects (TWFE) regression:
\begin{equation*}
Y_{i,t} = \alpha_{0}+\gamma_{0}\mathbf{1}\{G_i = A\}+\lambda_{0}\mathbf{1}\{T_i = 2023\}+\beta_{0}^{TWFE}\big(\mathbf{1}\{G_i = A\}\cdot\mathbf{1}\{T_i = 2023\}\big)+\varepsilon_{i,t}
\end{equation*}
\begin{enumerate}[(a)]
\item Interpret each parameter ($\alpha_{0}$, $\gamma_{0}$, $\lambda_{0}$, $\beta_{0}^{TWFE}$) in terms of conditional expectations of observed outcomes.
\item Show that under SUTVA, No-Anticipation, and Parallel Trends, $\beta_{0}^{TWFE}$ equals the ATT.
\item Explain why $\beta_{0}^{TWFE}$ equals the ``DiD by hand'' estimator you calculated in question 1(a).
\end{enumerate}

\end{enumerate}

\end{document}
